
\subsection{Model Initialization}

Before the advent of FSDP, PyTorch mandated the full materialization of the entire model instance on one device. Although users can allocate different sub-modules to different devices, this would require modifying the model source code, which may not be feasible, particularly if model authors and application developers belong to different parties. To facilitate a smooth transition from local to distributed training, FSDP must effectively aid in the materialization and initialization of a massive model, which poses two challenges:

\begin{itemize}
    \item How to create a model instance without materializing any tensor storage, postponing initialization until a storage on a concrete device is attached to the tensor.
    \item How to ensure accurate initialization of model parameters in line with the user's implementation, even when the model is too large to fit on a single GPU. 
\end{itemize}

To overcome the first challenge, we have introduced a mechanism called \emph{deferred initialization}, which involves the allocation of model parameter tensors on a simulated or "fake" device. During this process, all initialization operations performed on the tensor are recorded. Subsequently, when the tensor is moved from the "fake" device to a GPU device, all recorded operations are automatically replayed. By adopting this technique, users can generate a model instance from any third-party library without allocating any GPU memory blocks, while still accurately capturing their parameter initialization implementations.

As illustrated in Figure~\ref{fig:overview}, once the FSDP has wrapped the model, it is evenly distributed across all GPUs, with each device holding only one shard in its memory. Therefore, in order to address the second challenge, each rank should ideally only materialize and initialize the shard that it owns. However, this is not always practical, since we cannot predict what initialization logic the user will implement in the model \code{init} method. The initialization logic may rely on having a unsharded parameter on the device, which makes it impossible to shard the initialization. Consequently, FSDP must prepare the unsharded parameters before executing \code{Tensor} initialization operations and simultaneously reduce the memory footprint. Given that sharding initialization is unsafe, FSDP applies the same approach as how it handles model forward and backward passes, \emph{i.e.}, initialize one FSDP unit at a time and shard the unit before moving on to the next one. When combined with deferred initialization, FSDP traverses the fake device model instance to decompose it into FSDP units, moves one unit to a GPU device at a time, and replays the recorded initialization operations for tensors in that FSDP unit.



%%%%%%%%%%%%%


